\section{What can be universal?}
\label{sec:universal}

An explicit sparse vector with $k$ listed nonzeros requires $\Omega(k)$ output
words.  This elementary bound does not automatically become
$\Omega(\vol(S))$: a high-degree vertex still needs one output value unless
the requested certificate or cluster output requires scanning its incidences.
Every volume lower bound must state why adjacency exposure is necessary.

Similarly, no graph-uniform polynomial dependence on $1/\alpha$ should be
asserted merely from the condition number of $Q$.  Graph-aware SDD and
elimination algorithms use individual entries and response structure that a
black-box first-order oracle hides.  The 2026 growing-active-set SDD result also
rules out several uniform product lower bounds across all parameter regimes in
its model.

Meaningful universal lower bounds may instead depend on:

\begin{itemize}
\item explicit output size;
\item the number of adjacency incidences that must be distinguished;
\item allowed approximation and failure probability;
\item boundary-reporting output size;
\item memory, adaptivity rounds, or bit complexity;
\item a representation parameter such as response rank, separator size, or
      required refinement count.
\end{itemize}
